\chapter{Babel plugin in detail (JSX transform)}
\label{ch:babel-plugin}

This chapter explains the current HRBR Babel transform in \texttt{compiler/jsx/transform.ts}.
It is intentionally a \emph{scaffold}: it covers a useful subset, has crisp failure modes, and is wired end-to-end in jsdom tests.

The transform's job is to recognize TSX/JSX expressions that can be compiled into a stable \texttt{BlockDef} + \texttt{CompiledSlot[]} and to rewrite them into a \texttt{mount(host) => instance} factory that the VDOM interop layer can consume (see Chapter~\ref{ch:vdom-hrbr-interop}).

When it can't safely compile a JSX expression into a stable template, it routes to \texttt{mountFallback()}.

\section{Public entrypoint and options}

The plugin is exported as \texttt{hrbrJsxTransform}:

\begin{itemize}
\item \texttt{compiler/jsx/index.ts}: re-exports the default plugin
\item \texttt{compiler/index.ts}: re-exports it as part of the public compiler API
\end{itemize}

\subsection{Tiny contract: what this plugin must guarantee}

The transform is successful if it guarantees two things:

\begin{itemize}
\item The emitted program is \textbf{self-contained}: all runtime references are imported (or intentionally omitted when not needed).
\item Any transformed JSX expression becomes a \textbf{mount factory} of the form \texttt{(host: Element) => MountedInstance}, matching the HRBR vnode contract from Chapter~\ref{ch:vdom-hrbr-interop}.
\end{itemize}

Everything else (slot heuristics, fallback subset size, key readability) is an implementation choice that can evolve as long as the tests keep passing.

The options type \texttt{HrbrJsxTransformOptions} (in \texttt{compiler/jsx/transform.ts}) controls a few key behaviors:

\begin{itemize}
\item \texttt{runtimeImport}: module path for runtime helpers (default \texttt{"relax/hrbr"})
\item \texttt{mountWrapperName}: helper function name used to call \texttt{mountCompiledBlock} (default \texttt{\_\_hrbrBlock})
\item \texttt{dev}: enables more readable output and debug-oriented metadata
\item \texttt{stableSlotKeys}: when true, derive slot keys from source locations (defaults to \texttt{dev})
\end{itemize}

\section{High-level rewrite shape}

Given a JSX expression like:

\begin{quote}
\texttt{const App = () => <div className=\{cls()\}>Hi \{name()\}</div>}
\end{quote}

The plugin emits three conceptual things:

\begin{enumerate}
\item An import of runtime functions:\\
\texttt{import \{ defineBlock, mountCompiledBlock \} from "relax/hrbr"}
\item A module-scope \texttt{const \_\_hrbr\_block\_N = defineBlock (\ldots)} containing:
  \begin{itemize}
  \item \texttt{templateHTML: "<div ...>..."}
  \item \texttt{slots: \{ key: \{ kind, path, name?, lane? \}, ... \}}
  \end{itemize}
\item A replacement of the JSX expression with a mount factory:\\
\texttt{host => \_\_hrbrBlock(host, \_\_hrbr\_block\_N, [\{key, read\}, ...])}
\end{enumerate}

This shape matches the interop contract in \texttt{src/h.ts}: an HRBR vnode is just a mount function that returns an instance.

\subsection*{Diagram: one JSX expression becomes a hoisted block + mount factory}
\begin{center}
\begin{tikzpicture}[
  node distance=10mm and 14mm,
  box/.style={draw, rounded corners, align=left, inner sep=6pt, fill=RelaxLight},
  arrow/.style={-Latex, thick}
]
\node[box] (input) {Input code\\\texttt{const App = () => <div>Hi \{name()\}</div>}};
\node[box, right=of input] (visit) {Babel visitor\\\texttt{JSXElement(path)}};
\node[box, right=of visit] (emit) {Emit artifacts\\1) hoisted \texttt{defineBlock(...)}\\2) mount factory replaces JSX};

\node[box, below=of emit] (hoist) {Module scope\\\texttt{const \_\_hrbr\_block\_0 = defineBlock(...)};};
\node[box, below=of visit] (replace) {Expression position\\\texttt{(host) => \_\_hrbrBlock(host, \_\_hrbr\_block\_0, [\{key, read\}, ...])}};

\draw[arrow] (input) -- (visit);
\draw[arrow] (visit) -- (emit);
\draw[arrow] (emit) -- (hoist);
\draw[arrow] (emit) -- (replace);
\end{tikzpicture}
\end{center}

The key optimization is hoisting: the template is parsed once per module, while the mount factory runs per render/mount.

\section{Program-level state and synthesis (imports + wrappers)}

The plugin keeps per-file state variables:

\begin{itemize}
\item \texttt{needsRuntimeImport}: whether to inject \texttt{defineBlock/mountCompiledBlock}
\item \texttt{needsFallbackImport}: whether to also import \texttt{mountFallback}
\item \texttt{needsWrapper}: whether to emit the wrapper function \texttt{\_\_hrbrBlock}
\item \texttt{nextBlockId}: a counter used to generate \texttt{\_\_hrbr\_block\_0}, \texttt{\_\_hrbr\_block\_1}, ...
\item \texttt{currentFileName\/currentFileCode}: used for context-rich error messages
\end{itemize}

These are reset in the \texttt{Program.enter} visitor and materialized in \texttt{Program.exit}:

\begin{itemize}
\item \texttt{ensureRuntimeImport(programPath)} uses \texttt{unshiftContainer('body', ImportDeclaration)}.
\item \texttt{ensureWrapper(programPath)} appends \texttt{function \_\_hrbrBlock(host, def, slots) \{ return mountCompiledBlock(def, host, slots) \}}.
\end{itemize}

Implementation note: this scaffold avoids depending on Babel "types" helpers by emitting raw AST node objects.

\subsection*{Diagram: per-file state at \texttt{Program.enter/exit}}
\begin{center}
\begin{tikzpicture}[
  node distance=9mm and 12mm,
  box/.style={draw, rounded corners, align=left, inner sep=6pt, fill=RelaxLight},
  arrow/.style={-Latex, thick}
]
\node[box] (enter) {\texttt{Program.enter}\\reset file state:\\\texttt{needsRuntimeImport=false}\\\texttt{needsFallbackImport=false}\\\texttt{needsWrapper=false}\\\texttt{nextBlockId=0}};
\node[box, right=of enter] (walk) {Walk AST\\each \texttt{JSXElement}:\\\texttt{built = buildTemplateFromJsxElement(...)}\\toggle flags};
\node[box, right=of walk] (exit) {\texttt{Program.exit}\\materialize emissions:\\imports, wrapper helper,\\and \texttt{\_\_hrbr\_block\_N} constants};

\draw[arrow] (enter) -- (walk);
\draw[arrow] (walk) -- (exit);
\end{tikzpicture}
\end{center}

This structure prevents unused imports/helpers: nothing is emitted unless some subtree required it.

\subsection{Why inject at \texttt{Program.exit}?}

The plugin defers emissions until \texttt{Program.exit} so it can answer questions like:

\begin{itemize}
\item Did we ever successfully compile a block? If so, we need \texttt{defineBlock} and \texttt{mountCompiledBlock}.
\item Did any subtree route to fallback? If so, we need \texttt{mountFallback} too.
\item Did we generate a wrapper call? If so, we need to emit the wrapper function body.
\end{itemize}

This design keeps the output minimal: files that contain no transformable JSX don't get imports they don't use.

\section{The core visitor: \texttt{JSXElement(path)}}

The plugin's heart is the \texttt{JSXElement} visitor.
For each JSX element in an expression position, it calls:

\begin{quote}
\texttt{const built = buildTemplateFromJsxElement(path)}
\end{quote}

\subsection{Success path: compile into a block}

When \texttt{built} is non-null:

\begin{enumerate}
\item Mark \texttt{needsRuntimeImport = true} and \texttt{needsWrapper = true}.
\item Insert a module-scope \texttt{const \_\_hrbr\_block\_N = defineBlock(\{ templateHTML, slots \})}.
\item Replace the JSX expression with:\\
\texttt{host => \_\_hrbrBlock(host, \_\_hrbr\_block\_N, built.compiledSlots)}
\end{enumerate}

The slot objects are turned into AST literals through \texttt{literalizeSlot(slot)}.

\subsubsection{Hoisting: why \texttt{defineBlock(...)} lives at module scope}

The plugin uses \texttt{unshiftContainer('body', ...)} on the program node to emit:

\begin{quote}
	exttt{const \_\_hrbr\_block\_N = defineBlock({ templateHTML, slots })}
\end{quote}

Hoisting matters for two reasons:

\begin{itemize}
\item \textbf{Performance:} we want template parsing and slot-path validation to happen once per module, not once per render.
\item \textbf{Identity:} the mount factory will refer to a stable \texttt{\_\_hrbr\_block\_N} constant.
\end{itemize}

\subsection{Failure path: route to fallback}

When \texttt{built} is null, the plugin attempts a \emph{minimal fallback transform}.

\paragraph{Why fallback exists}

Block mode requires a stable template. Expressions that can produce different node shapes at runtime (conditionals, JSX inside expressions, etc.) must be executed by the fallback runtime, which uses the reconciler.

\subsection*{Diagram: fallback routing when template extraction fails}
\begin{center}
\begin{tikzpicture}[
  node distance=10mm and 14mm,
  box/.style={draw, rounded corners, align=left, inner sep=6pt, fill=RelaxLight},
  arrow/.style={-Latex, thick},
  decision/.style={draw, diamond, aspect=2, inner sep=2.5pt, align=center, fill=RelaxLight}
]
\node[box] (jsx) {JSX subtree\\\texttt{<div>...}};
\node[decision, right=of jsx] (ok) {Can build\\stable template?};
\node[box, above right=of ok] (block) {Block transform\\emit \texttt{defineBlock + compiledSlots}};
\node[box, below right=of ok] (fb) {Fallback transform\\emit \texttt{host => mountFallback(...)}\\\texttt{renderFn} returns specs};

\draw[arrow] (jsx) -- (ok);
\draw[arrow] (ok) -- node[midway, above] {yes} (block);
\draw[arrow] (ok) -- node[midway, below] {no} (fb);
\end{tikzpicture}
\end{center}

In other words: block mode is an optimization; fallback mode is the semantic firewall.

\paragraph{Current fallback subset}

The fallback routing in \texttt{transform.ts} is intentionally small:

\begin{itemize}
\item It only supports intrinsic lowercase root tags.
\item It only supports children that are \texttt{JSXText} or \texttt{JSXExpressionContainer}.
\item It allows only static string attributes (and ignores dynamic attrs).
\item It rejects spread attributes and event handlers.
\end{itemize}

In this subset, the plugin synthesizes a \texttt{render()} function (a closure-free data + expression composition) and emits:

\begin{quote}
\texttt{host => mountFallback(host, () => (\{ children: [spec] \}))}
\end{quote}

Each child spec contains:

\begin{itemize}
\item \texttt{create(): Element} -> creates the element and sets attributes/text
\item \texttt{patch(node)} -> updates textContent based on the latest expressions
\item optional \texttt{key} if a \texttt{key=} attribute is present
\end{itemize}

This is enough to prove the end-to-end pipeline and to connect the compiler to the fallback runtime from earlier chapters.

\subsection{Fallback synthesis: what gets captured (and what doesn't)}

The fallback compiler intentionally avoids capturing complex structures.
It builds a \texttt{renderFn} with the shape:\\
	exttt{() => ({ children: [spec] })}

Where \texttt{spec.create()} and \texttt{spec.patch(node)} both close over the user expressions.

Important implications:

\begin{itemize}
\item \textbf{Reactivity comes from runtime:} \texttt{mountFallback} re-runs \texttt{renderFn} under tracking and schedules reconciliation.
\item \textbf{This is not a general JSX renderer:} it is a deliberate minimal bridge used to validate the ``block or fallback'' architecture.
\end{itemize}

\section{Template extraction: \texttt{buildTemplateFromJsxElement()}}

The function \texttt{buildTemplateFromJsxElement()} is a best-effort compiler for a supported subset.
It returns either a full block build or \texttt{null}.

\subsection{Supported subset (today)}

In block mode, the transformer supports:

\begin{itemize}
\item intrinsic lowercase tags only (e.g. \texttt{div}, \texttt{section})
\item nested elements
\item static string attributes
\item expression attributes -> slot kinds \texttt{class/style/attr}
\item dynamic text children -> only when the expression is ``text-like''
\end{itemize}

It explicitly rejects (in block mode):

\begin{itemize}
\item event handlers (\texttt{onClick=} etc.)
\item fragments (\texttt{<>...</>})
\item non-intrinsic tags (components, capitalized)
\item spread attributes (\texttt{\{...props\}})
\end{itemize}

\subsection{How paths are computed in the Babel transform}

Like the TS prototype compiler (Chapter~\ref{ch:compiler-architecture}), this transform must produce \texttt{childNodes}-based paths.

In \texttt{renderElement(...)} the transform:

\begin{itemize}
\item keeps a per-parent counter \texttt{childState.i} that increments for each child node
\item uses \texttt{pathToHere} plus the current index to form paths like \texttt{[3,1]}
\item emits a placeholder text node (a single space) when it needs to guarantee the existence of a text node for a dynamic text slot
\end{itemize}

This is why the snapshot test shows a template that contains \emph{extra whitespace} in places where text slots are inserted. That whitespace isn't an aesthetic choice; it's a structural anchor.

\subsection{Slot keys: sequential vs location-based}

Slot keys are allocated in \texttt{allocSlotKey()}.

\begin{itemize}
\item Default: sequential keys \texttt{s0}, \texttt{s1}, ...
\item Dev mode (\texttt{stableSlotKeys=true}): keys include source location: \texttt{s\_\textless{}kind\textgreater{}\_\textless{}line\textgreater:\textless{}col\textgreater}.
\end{itemize}

This is validated in \texttt{compiler/\_\_tests\_\_/babel-transform.test.ts}.

One nuance: location-based slot keys depend on Babel preserving \texttt{loc} information and comments; this is why dev mode is framed as "debug output" rather than a semantic requirement.

\subsection{Lane pragmas: \texttt{/*@lane transition*/}}

The helper \texttt{extractLanePragmaFromExpressionContainer()} scans leading comments and looks for:

\begin{quote}
\texttt{@lane (sync|input|default|transition|idle)}
\end{quote}

If present, the emitted slot metadata includes \texttt{lane: "..."}.
The tests verify that the transformed output contains the lane strings.

\subsection{Text-like safety}

In \texttt{isTextLikeExpr()}, the transform tries to reject expressions that can yield JSX elements or otherwise change DOM structure.

Examples treated as \emph{not text-like}:

\begin{itemize}
\item returning JSX elements/fragments
\item conditionals/logicals that could yield elements (conservatively rejected)
\end{itemize}

When an expression is accepted as text-like, the transform:

\begin{itemize}
\item emits a sentinel space in the template (so a text node exists)
\item records a \texttt{text} slot pointing to that position
\item emits a compiled slot \texttt{read: () => expr}
\end{itemize}

The predicate \texttt{isTextLikeExpr(expr)} is the key safety gate.
It explicitly rejects:

\begin{itemize}
\item \texttt{JSXElement} and \texttt{JSXFragment} (obviously structural)
\item \texttt{ConditionalExpression} and \texttt{LogicalExpression} (conservatively structural for now)
\end{itemize}

This conservative choice is why the fallback e2e test uses \texttt{show() ? name() : ''} to force fallback routing: it is not considered "text-like" yet.

\section{Failure modes are part of the API}

This compiler phase aims for clarity and determinism.
Unsupported syntax produces explicit errors via \texttt{fail(path, message)} which attempts to build Babel code-frame errors.

The file \texttt{compiler/\_\_tests\_\_/babel-transform.errors.test.ts} locks in these failure modes:

\begin{itemize}
\item event handlers should throw
\item spread attributes should throw
\item fragments should throw
\item non-intrinsic tags should throw
\end{itemize}

While these errors might eventually become "route to fallback" behaviors, they're still valuable today: they prevent the compiler from silently generating templates that would be incorrect (for example, by treating an event handler like an attribute).

\section{End-to-end tests: jsdom execution proves output shape}

Two e2e suites validate that the emitted output is not just syntactically correct but \emph{runnable}.

\subsection{Block mode e2e: transformed output mounts and updates}

\texttt{compiler/\_\_tests\_\_/babel-transform.e2e-jsdom.test.ts} compiles a small component, extracts the generated \texttt{defineBlock(...)} literal and the \texttt{compiledSlots} array literal, evaluates them with \texttt{new Function(...)} and mounts the result with \texttt{mountCompiledBlock}.

The key invariant: updating a signal updates the DOM without recreating the root element.

This suite also proves a deeper property: the emitted output can be \emph{partially evaluated}.
The test extracts only the \texttt{defineBlock(...)} object literal and the compiled slots array literal, then evaluates them with \texttt{new Function}.
So the plugin isn't just emitting runnable code; it's emitting code whose pieces behave like data.

\subsection{Fallback routing e2e: structural text expression goes to \texttt{mountFallback}}

\texttt{compiler/\_\_tests\_\_/babel-transform.fallback.e2e-jsdom.test.ts} compiles an expression with a structural conditional:

\begin{quote}
\texttt{<div>\{show() ? name() : ''\}</div>}
\end{quote}

This should include \texttt{mountFallback} in output and remain reactive under signal updates.

The spec here is intentionally small:

\begin{itemize}
\item it asserts that \texttt{mountCompiledBlock} should \emph{not} be called
\item it asserts that mounting succeeds and subsequent signal changes cause updates
\item it does \emph{not} assert perfect minimal DOM operations yet
\end{itemize}

That matches the current purpose of fallback in this compiler phase: correctness-first safety for dynamic structure.

\section{What to improve next (guided by tests)}

The current transform is intentionally conservative.
The next incremental steps typically look like:

\begin{itemize}
\item add safe event slot support (mapping \texttt{onClick} to HRBR event slots)
\item support fragments by lowering them to concatenated children in block mode
\item expand fallback routing to support nested children specs (not just a single root spec)
\item improve text-like analysis so common conditionals can stay in block mode when both arms are text
\end{itemize}

The rule is always: add a targeted test first (unit + e2e), then widen syntax support.

When you widen the subset, remember the three invariants all previous chapters rely on:

\begin{itemize}
\item slot paths must match \texttt{childNodes}
\item stable mount function identity enables reuse in VDOM interop
\item fallback must remain safe when structure can change
\end{itemize}

\section{Online resources}

\begin{itemize}
\item Babel plugin authoring and visitor lifecycle:\\
\href{https://babeljs.io/docs/en/plugins}{Babel docs: plugins},\\
\href{https://babeljs.io/docs/en/babel-plugin-api}{Babel docs: plugin API}.

\item JSX AST node shapes (useful when reading \texttt{JSXElement}, \texttt{JSXText}, and \texttt{JSXExpressionContainer} visitor logic):\\
\href{https://github.com/babel/babel/blob/main/packages/babel-parser/ast/spec.md}{Babel parser AST spec}.

\item Source maps and why transforms try to preserve locations:\\
\href{https://sourcemaps.info/spec.html}{Source Map Revision 3 Proposal}.

\item ``Static template + dynamic slots" as a general pattern (useful for intuition):\\
\href{https://svelte.dev/docs/svelte/what-is-svelte}{Svelte docs: compiler-first model}.
\end{itemize}
